% RACER-GT 使用手冊（繁體中文）。
\documentclass[12pt,a4paper,oneside]{article}
% Traditional Chinese setup. Compile with XeLaTeX.
\usepackage{fontspec}
\usepackage{xeCJK}
\setmainfont{Noto Serif}
\setsansfont{Noto Sans}
\setmonofont{Noto Sans Mono}
\setCJKmainfont{Noto Serif CJK TC}
\setCJKsansfont{Noto Sans CJK TC}
\setCJKmonofont{Noto Sans CJK TC}
\XeTeXlinebreaklocale "zh"
\XeTeXlinebreakskip = 0pt plus 1pt

% Single source of truth for version and date across every RACER-GT document.
% Regenerate nothing by hand: bump here, rebuild, and every PDF agrees.
\newcommand{\racerversion}{1.1.0}
\newcommand{\racerdate}{2026-07-27}
\newcommand{\racerauthor}{Yi-Hao Lai}
\newcommand{\raceraffil}{Department of Finance, Da-Yeh University}
\newcommand{\raceremail}{yhlai@mail.dyu.edu.tw}
\newcommand{\racerrepo}{https://github.com/yhlai0911/RACER-GT}

% Shared package set and mathematical macros for every RACER-GT document,
% in both languages. Language-specific font and theorem-name setup lives in
% _preamble-zh.tex and _preamble-en.tex.
%
% Font policy: the build depends only on the five Noto families below. They are
% installable on Linux CI (fonts-noto-core, fonts-noto-cjk) and on macOS via
% Homebrew casks. Do not introduce a font outside this set without adding it to
% .github/workflows/docs.yml as well, or the PDF build stops being reproducible.

\usepackage[margin=2.35cm,headheight=15pt]{geometry}
\usepackage{amsmath,amssymb,amsthm,mathtools,bm}
\usepackage{booktabs,longtable,tabularx,array,multirow}
\usepackage{graphicx,float,caption,subcaption}
\usepackage{enumitem}
\usepackage{xcolor}
\usepackage{microtype}
\usepackage{setspace}
\usepackage{fancyhdr}
\usepackage{hyperref}
\usepackage[nameinlink,noabbrev]{cleveref}
\usepackage{algorithm}
\usepackage{algpseudocode}
\usepackage{listings}
\usepackage{tikz}
\usetikzlibrary{arrows.meta,positioning,shapes.geometric,fit,calc}

\hypersetup{
  colorlinks=true,
  linkcolor=black,
  citecolor=black,
  urlcolor=blue,
  pdfauthor={\racerauthor}
}

\pagestyle{fancy}
\fancyhf{}
\rhead{v\racerversion}
\cfoot{\thepage}
\setstretch{1.28}
\setlength{\parskip}{0.35em}
\setlist[itemize]{leftmargin=2em,itemsep=0.25em,topsep=0.3em}
\setlist[enumerate]{leftmargin=2.2em,itemsep=0.25em,topsep=0.3em}

% ---------------------------------------------------------------- mathematics
\newcommand{\E}{\mathbb{E}}
\newcommand{\Var}{\operatorname{Var}}
\newcommand{\Cov}{\operatorname{Cov}}
\newcommand{\tr}{\operatorname{tr}}
\newcommand{\diag}{\operatorname{diag}}
\newcommand{\argmin}{\operatorname*{arg\,min}}
\newcommand{\argmax}{\operatorname*{arg\,max}}
\newcommand{\one}{\bm{1}}
\newcommand{\Rset}{\mathbb{R}}
\newcommand{\R}{\mathbb{R}}
\newcommand{\plim}{\operatorname*{plim}}
\newcommand{\RACER}{\textsc{Racer-GT}}
\newcommand{\indic}[1]{\mathbb{I}\!\left\{#1\right\}}
\newcommand{\Sig}{\bm{\Sigma}}
\newcommand{\wvec}{\bm{w}}
\newcommand{\evec}{\bm{e}}
\newcommand{\Zvec}{\bm{Z}}

% Design facets, used identically in both languages so that a reader can move
% between the Chinese and English documents without re-learning notation.
\newcommand{\facetT}{\mathcal{T}}   % historical date
\newcommand{\facetD}{\mathcal{D}}   % collection day
\newcommand{\facetS}{\mathcal{S}}   % collection stream
\newcommand{\facetR}{\mathcal{R}}   % technical replicate

\lstset{
  basicstyle=\ttfamily\footnotesize,
  breaklines=true,
  frame=single,
  rulecolor=\color{gray!40},
  backgroundcolor=\color{gray!5},
  columns=fullflexible,
  keepspaces=true,
  showstringspaces=false
}


\setlength{\parindent}{2em}

\newtheorem{assumption}{假設}[section]
\newtheorem{proposition}{命題}[section]
\newtheorem{corollary}{推論}[section]
\newtheorem{remark}{說明}[section]
\newtheorem{definition}{定義}[section]
\newtheorem{lemma}{引理}[section]

\renewcommand{\proofname}{證明}
\renewcommand{\abstractname}{摘要}
\renewcommand{\contentsname}{目錄}
\renewcommand{\refname}{參考文獻}
\renewcommand{\tablename}{表}
\renewcommand{\figurename}{圖}

\crefname{assumption}{假設}{假設}
\crefname{proposition}{命題}{命題}
\crefname{corollary}{推論}{推論}
\crefname{definition}{定義}{定義}
\crefname{lemma}{引理}{引理}
\crefname{equation}{式}{式}
\crefname{table}{表}{表}
\crefname{figure}{圖}{圖}
\crefname{section}{節}{節}

\lhead{RACER-GT：使用手冊}

\title{\bfseries RACER-GT 使用手冊\\[0.4em]
\large 安裝、工作流程、資料契約與輸出}
\author{\racerauthor\\\small 大葉大學 財務金融學系}
\date{\racerdate\ \textbullet\ 版本 \racerversion}

\begin{document}
\maketitle

\begin{abstract}
\noindent 本手冊帶研究者從安裝走到一條通過驗收的日頻 Google Trends 指標。內容涵蓋收集器必須滿足的資料契約、四個命令列步驟、Python API、每一個輸出檔案，以及實務上最容易遇到的失效情況。方法與證明請見方法論主文與數學附錄；本文只談如何正確操作軟體。
\end{abstract}

\tableofcontents
\newpage

% =============================================================================
\section{安裝}

RACER-GT 需要 Python 3.10 以上。

\begin{lstlisting}[language=bash]
# 由 PyPI 安裝
python -m pip install racergt

# 由 release wheel 安裝
python -m pip install racergt-1.0.0-py3-none-any.whl

# 由原始碼安裝（開發用）
python -m pip install -e '.[dev]'
pytest -q
\end{lstlisting}

本套件刻意採 \emph{ingestion-first}：不內含任何 Google Trends 抓取程式。端點、rate limit、認證與服務條款都可能改變，統計方法若綁定非官方 client，就會繼承那份脆弱性。收集器由您選擇；套件只要求它輸出 \cref{sec:contract-zh} 的 schema。

% =============================================================================
\section{四步驟工作流程}

\subsection{步驟一：鎖定 protocol}

\begin{lstlisting}[language=bash]
racergt design config.yaml --out protocol_bundle --anchor-date 2026-08-01
\end{lstlisting}

產出帶 SHA-256 protocol hash 的 \texttt{protocol.lock.yaml}、平衡後的 day $\times$ stream $\times$ replicate 計畫 \texttt{collection\_schedule.csv}、固定查詢視窗 \texttt{chunk\_windows.csv}，以及稽核 manifest。

\begin{remark}[不要手動編輯 lock 檔]
\texttt{RacerGTConfig.load\_yaml} 會重算 hash 並在不符時拋錯。手動改動被鎖定的設定因此會明確失敗，這是刻意的行為。要改設定，請改來源設定檔後重新產生。
\end{remark}

\subsection{步驟二：蒐集與轉換}

依 schedule 執行。每個 response 原封不動保存，再轉為 \cref{sec:contract-zh} 的 long format。每個 chunk 都計算原始 response 與數值向量的 hash。若發生網路中斷、登入改變、cookie 清除或 client 更新，請記錄於 metadata；不要自行判定該 pull 無效——那會構成依賴數值的排除。

\subsection{步驟三：稽核}

\begin{lstlisting}[language=bash]
racergt audit config.yaml raw_chunks.csv --out audit.json
\end{lstlisting}

稽核在任何估計之前，先比對批次與已鎖定 protocol。失敗時 exit code 為 2。

\subsection{步驟四：執行}

\begin{lstlisting}[language=bash]
racergt run config.yaml raw_chunks.csv \
  --benchmark lower_frequency_benchmark.csv \
  --out results
\end{lstlisting}

若鎖定的 decision tree 回傳 FAIL，exit code 為 3，可直接用於自動化流程的閘門。

\subsection{Python API}

\begin{lstlisting}[language=Python]
import pandas as pd
from racergt import RacerGTConfig, RacerGTPipeline

config = RacerGTConfig.load_yaml("config.yaml")
raw = pd.read_csv("raw_chunks.csv")
benchmark = pd.read_csv("lower_frequency_benchmark.csv")

result = RacerGTPipeline(config).fit(raw, benchmark=benchmark)
result.save("results")

print(result.decision.status)      # PASS / REVIEW / FAIL
print(result.final_series.head())
\end{lstlisting}

% =============================================================================
\section{資料契約}\label{sec:contract-zh}

\subsection{Raw chunk table}

每列對應一個 \texttt{pull\_id} $\times$ \texttt{chunk\_id} $\times$ \texttt{historical\_date}。

\begin{table}[H]
\centering\small
\caption{必要欄位。}
\begin{tabularx}{\textwidth}{llX}
\toprule
欄位 & 型別 & 意義\\
\midrule
\texttt{series\_id} & string & 研究指標識別碼\\
\texttt{pull\_id} & string & 一次完整歷史重建\\
\texttt{chunk\_id} & string & 固定查詢視窗，須存在於 manifest\\
\texttt{historical\_date} & date & GT 觀察日期\\
\texttt{value} & float & 原始 GT 相對搜尋值，通常 0--100\\
\texttt{collection\_day} & integer & 蒐集日的設計層級\\
\texttt{stream\_id} & string & 固定的複合蒐集環境\\
\texttt{replicate\_id} & string & day $\times$ stream 格內的 technical replicate\\
\texttt{window\_start} & date & chunk 起日\\
\texttt{window\_end} & date & chunk 迄日\\
\bottomrule
\end{tabularx}
\end{table}

\begin{remark}[兩個經常被誤解的欄位]
\texttt{collection\_day} 是蒐集日的\emph{序數設計層級}（0、1、2……），\textbf{不是} day offset，也不是日曆日期。\texttt{stream\_id} 代表複合環境——裝置、瀏覽器 profile、cookie、帳號狀態、網路路徑、IP。除非確實執行了 crossover factorial design，否則估出的 stream effect 不得報告為 IP effect。
\end{remark}

強烈建議另存：\texttt{retrieved\_at}、\texttt{protocol\_hash}、\texttt{keyword}、\texttt{geo}、\texttt{category}、\texttt{search\_property}、\texttt{language}、\texttt{is\_partial}、\texttt{raw\_response\_hash}、\texttt{network\_event}。

\subsection{低頻 benchmark table}

需要 \texttt{series\_id}、\texttt{period\_start}、\texttt{period\_end}、\texttt{value}；\texttt{se} 為選填，缺省時採 protocol 預設值。

% =============================================================================
\section{輸出}

\begin{table}[H]
\centering\small
\caption{\texttt{result.save()} 產生的檔案。}
\begin{tabularx}{\textwidth}{lX}
\toprule
路徑 & 內容\\
\midrule
\texttt{RACER\_GT\_final\_series.csv} & 最終日頻指標、條件式標準誤、95\% 區間\\
\texttt{complete\_calibrated\_pulls.csv} & 全部校準後的 pull\\
\texttt{calibration/*\_chunk\_scales.csv} & 各 chunk 的全域尺度估計\\
\texttt{calibration/*\_overlap\_edges.csv} & 穩健 edge 估計與診斷\\
\texttt{duplicates/} & exact groups、pair metrics、components、殘差\\
\texttt{consensus/design\_weighted\_consensus.csv} & 設計式參考估計量\\
\texttt{consensus/consensus\_weights.csv} & GLS／最小變異權重\\
\texttt{gstudy/} & 各轉換的 ANOVA、變異成分、D-study\\
\texttt{reliability/} & pairwise reliability、收斂路徑、day/stream 相依\\
\texttt{decision/acceptance\_decision\_tree.csv} & 每條規則的觀察值與門檻\\
\texttt{report/RACER\_GT\_diagnostic\_report.md} & 圖表式診斷報告\\
\texttt{summary.json} & protocol hash、狀態、主要診斷\\
\bottomrule
\end{tabularx}
\end{table}

\begin{remark}[兩個 consensus 檔案，但只有一條最終數列]
\texttt{consensus/design\_weighted\_consensus.csv} 是設計式參考估計量，僅供對照，\textbf{不會}進入 \texttt{RACER\_GT\_final\_series.csv}。後者在提供 benchmark 時來自 temporal benchmark，否則來自 GLS consensus。兩者若差距很大，是值得追查的警訊，不是程式錯誤。
\end{remark}

% =============================================================================
\section{解讀 decision}

任一 mandatory 規則失敗即 FAIL；無失敗但有 indeterminate 則 REVIEW；否則 PASS。常見成因與處置：

\begin{table}[H]
\centering\small
\caption{常見結果與對應處置。}
\begin{tabularx}{\textwidth}{lX}
\toprule
現象 & 解讀與處置\\
\midrule
\texttt{OverlapGraphError} & chunk 設計有缺口。加長視窗或縮短 step，使相鄰 chunk 至少共享 \texttt{min\_overlap\_days} 個可用日。\\
Spectral 有效 pull 數過低 & 重複取得高度相依。應增加 collection day 或 stream；增加 technical replicate 無濟於事。\\
Detection 規則 indeterminate & 無零／非零變異，$\kappa$ 無定義。此情況下規則自動轉為非強制，無需處理。\\
零值占比過高 & 該詞彙在日頻下太稀疏。改用週頻，或改查較廣的 Topic 而非 Term。\\
Innovation G 未達標但 level G 達標 & 水準可重現、日變動不可重現。不要以差分形式使用該數列。\\
Benchmark RMSE 過大 & 日頻 consensus 與低頻數列不一致。請確認兩者使用完全相同的查詢設定。\\
\bottomrule
\end{tabularx}
\end{table}

% =============================================================================
\section{Stata 互通}

所有表格皆可輸出為 Stata 118 格式：

\begin{lstlisting}[language=bash]
racergt export-stata results/RACER_GT_final_series.csv final_series.dta
\end{lstlisting}

可讀格式：CSV、Parquet、XLS/XLSX；可寫格式：CSV、Parquet、XLSX、Stata 118 \texttt{.dta}。

% =============================================================================
\section{研究執行上的必要限制}

以下並非風格偏好；違反這些條件會使本框架提供的保證失效。

\begin{itemize}
\item 固定 IP 不構成獨立樣本的證據。
\item 不得在看過財金結果後調整 reliability 門檻或挑選 pull。
\item 若 pilot 監控導致 protocol 或 decision rule 改變，被監控的資料即成為 pilot 資料，正式批次必須重新開始。
\item exact duplicates 必須保留於 audit archive。
\item PASS 僅代表該批資料通過已鎖定的量測規則，不代表關鍵字具 construct validity，也不允許任何因果宣稱。
\end{itemize}

\end{document}
